% ==============================================================================
% Formation C# & SQL — Module 001 : Les Bases
% Fichier : src/P1_13_a.tex
% Sujet   : Chapitre 13 : Les Génériques
% ==============================================================================

\newpage
\section{Les Génériques}
\marginpar{\tiny\color{gray}P1\_13\_a.tex}

\subsection{Introduction Théorique}

Les génériques (\textit{Generics}) constituent un mécanisme fondamental du système de types C\# permettant de définir des classes, des structures, des interfaces et des méthodes dont un ou plusieurs types sont \textbf{paramétrés}. Au lieu de figer un type concret dans le code, le développeur déclare un espace réservé (conventionnellement noté \texttt{T}) qui sera substitué par un type réel lors de l'utilisation.

Ce mécanisme élimine deux problèmes majeurs de l'ingénierie logicielle :
\begin{itemize}
    \item \textbf{La duplication de code} : Sans génériques, il faudrait écrire une version de chaque algorithme pour chaque type de données (une méthode de tri pour \texttt{int}, une autre pour \texttt{string}, une autre pour \texttt{decimal}).
    \item \textbf{Le boxing/unboxing} : Avant les génériques (C\# 1.0), les collections comme \texttt{ArrayList} stockaient des \texttt{object}, forçant des conversions coûteuses en performance CPU entre types valeur et types référence.
\end{itemize}

Les collections que vous avez déjà utilisées (\texttt{List<T>}, \texttt{Dictionary<TKey, TValue>}) sont elles-mêmes des classes génériques fournies par le framework .NET.

\subsection{Syntaxe et Règles}

\subsubsection{Méthodes Génériques}

Une méthode générique déclare un paramètre de type entre chevrons (\texttt{<T>}) juste après son nom. Ce paramètre devient un type utilisable dans la signature et le corps de la méthode. Le compilateur infère automatiquement le type concret à partir des arguments passés à l'appel.

\subsubsection{Classes et Structures Génériques}

Une classe générique déclare ses paramètres de type au niveau de la classe elle-même. Toutes les méthodes et propriétés internes peuvent alors exploiter ces paramètres de type sans les redéclarer.

\subsubsection{Contraintes de Type (\texttt{where})}

Par défaut, un paramètre de type \texttt{T} peut être substitué par n'importe quel type existant, ce qui limite les opérations possibles sur \texttt{T} (on ne peut appeler que les méthodes de \texttt{object}). Les \textbf{contraintes} restreignent les types autorisés et débloquent des fonctionnalités supplémentaires :

\begin{center}
    \textbf{Tableau : Contraintes de type les plus courantes} \\[0.3cm]
    \begin{tabular}{@{}lp{9cm}@{}}
        \toprule
        \textbf{Contrainte} & \textbf{Effet} \\
        \midrule
        \texttt{where T : class} & \texttt{T} doit être un type référence (classe, interface, delegate). \\
        \texttt{where T : struct} & \texttt{T} doit être un type valeur non nullable. \\
        \texttt{where T : new()} & \texttt{T} doit posséder un constructeur public sans paramètre. Permet d'écrire \texttt{new T()} dans le code. \\
        \texttt{where T : IComparable<T>} & \texttt{T} doit implémenter l'interface spécifiée, donnant accès à ses méthodes. \\
        \texttt{where T : BaseClass} & \texttt{T} doit hériter de la classe de base spécifiée. \\
        \texttt{where T : notnull} & \texttt{T} ne peut pas être \texttt{null} (type valeur ou type référence non nullable). \\
        \bottomrule
    \end{tabular}
\end{center}

\subsection{Exemple de Code Commenté}

L'exemple suivant implémente un composant de cache générique en mémoire avec politique d'expiration, démontrant les méthodes génériques, les classes génériques et les contraintes de type.

\begin{lstlisting}[caption={Méthode générique et classe générique avec contraintes}]
// 1. Méthode générique simple : Recherche du maximum dans un tableau
// Contrainte : T doit implémenter IComparable<T> pour pouvoir être comparé
public static T TrouverMaximum<T>(T[] valeurs) where T : IComparable<T>
{
    if (valeurs.Length == 0)
        throw new ArgumentException("Le tableau ne peut pas être vide.", nameof(valeurs));

    T max = valeurs[0];
    for (int i = 1; i < valeurs.Length; i++)
    {
        // CompareTo() est accessible grâce à la contrainte IComparable<T>
        if (valeurs[i].CompareTo(max) > 0)
            max = valeurs[i];
    }
    return max;
}

// Appel : le compilateur infère T = int automatiquement
int[] mesures = { 42, 17, 93, 8, 65 };
int plusGrand = TrouverMaximum(mesures); // 93

// Appel explicite avec un autre type
string[] serveurs = { "alpha", "zulu", "bravo" };
string dernier = TrouverMaximum<string>(serveurs); // "zulu" (ordre lexicographique)
\end{lstlisting}

\begin{lstlisting}[caption={Classe générique : Cache en mémoire avec expiration}]
// 2. Classe générique avec deux paramètres de type
// Contrainte : TKey doit être non null (clé de cache), TValue doit avoir un constructeur
public class MemoryCache<TKey, TValue> 
    where TKey : notnull 
    where TValue : new()
{
    // Le dictionnaire interne utilise les paramètres de type de la classe
    private readonly Dictionary<TKey, (TValue Data, DateTime Expiration)> _store = new();

    // 3. Méthode utilisant les paramètres de type de la classe
    public void Set(TKey key, TValue value, TimeSpan ttl)
    {
        _store[key] = (value, DateTime.UtcNow.Add(ttl));
    }

    // 4. Retourne la valeur ou une nouvelle instance par défaut (grâce à new())
    public TValue GetOrDefault(TKey key)
    {
        if (_store.TryGetValue(key, out var entry) && entry.Expiration > DateTime.UtcNow)
        {
            return entry.Data;
        }
        return new TValue(); // Possible uniquement grâce à la contrainte 'where TValue : new()'
    }

    public int Count => _store.Count;
}

class Program
{
    static void Main()
    {
        // 5. Instanciation : substitution de TKey par string, TValue par List<string>
        var cache = new MemoryCache<string, List<string>>();

        var permissions = new List<string> { "READ", "WRITE", "ADMIN" };
        cache.Set("user_42_perms", permissions, TimeSpan.FromMinutes(30));

        List<string> perms = cache.GetOrDefault("user_42_perms");
        Console.WriteLine($"Permissions en cache : {string.Join(", ", perms)}");

        // 6. Tentative avec une clé expirée ou inexistante
        List<string> inconnu = cache.GetOrDefault("user_99_perms");
        Console.WriteLine($"Résultat par défaut : {inconnu.Count} éléments"); // 0 (new List<string>())
    }
}
\end{lstlisting}

\textbf{Explications ligne par ligne :}
\begin{itemize}
    \item \textbf{Ligne 3} : La méthode \texttt{TrouverMaximum<T>} déclare un paramètre de type \texttt{T} contraint par \texttt{IComparable<T>}. Sans cette contrainte, l'appel à \texttt{CompareTo()} provoquerait une erreur de compilation car le compilateur ne saurait pas que \texttt{T} possède cette méthode.
    \item \textbf{Ligne 19} : Le compilateur C\# infère automatiquement \texttt{T = int} à partir du type du tableau passé en argument. L'écriture explicite \texttt{TrouverMaximum<int>(mesures)} est équivalente mais redondante.
    \item \textbf{Lignes 28-30} : La classe \texttt{MemoryCache} déclare deux paramètres de type avec des contraintes distinctes. \texttt{notnull} garantit que la clé de cache ne sera jamais nulle. \texttt{new()} autorise l'instanciation d'un \texttt{TValue} par défaut.
    \item \textbf{Ligne 35} : Le dictionnaire interne utilise un tuple nommé \texttt{(TValue Data, DateTime Expiration)} pour stocker la valeur et sa date d'expiration. Les paramètres de type \texttt{TKey} et \texttt{TValue} de la classe sont directement exploitables dans ses membres.
    \item \textbf{Ligne 43} : L'instruction \texttt{new TValue()} est uniquement possible grâce à la contrainte \texttt{where TValue : new()}. Sans cette contrainte, le compilateur refuserait l'instanciation car il ne peut pas garantir que \texttt{TValue} possède un constructeur sans paramètre.
\end{itemize}

\begin{tcolorbox}[colback=white,colframe=keywordcolor,title=\textbf{Pièges courants \& Erreurs de compilation}]
    \begin{itemize}
        \item \textbf{\texttt{CS1061} — Membre inaccessible sur un type non contraint} : Tenter d'appeler \texttt{CompareTo()} ou toute méthode spécifique à une interface sur un \texttt{T} non contraint. Le compilateur restreint l'accès aux seuls membres hérités de \texttt{object} (\texttt{ToString()}, \texttt{Equals()}, \texttt{GetHashCode()}).
        \item \textbf{\texttt{CS0314} — Violation de contrainte au site d'appel} : Transmettre un argument de type qui ne satisfait pas la contrainte requise par une classe ou méthode générique.
        \item \textbf{\texttt{CS0310} — \texttt{new()} impossible} : Écrire \texttt{new T()} sans avoir déclaré la contrainte \texttt{where T : new()} dans la signature.
        \item \textbf{Confusion \texttt{class} vs \texttt{struct}} : Appliquer \texttt{where T : class} empêche l'utilisation de types valeur (\texttt{int}, \texttt{decimal}, \texttt{struct} personnalisées). Inversement, \texttt{where T : struct} interdit les types référence. Choisissez la contrainte en fonction de votre besoin d'allocation mémoire.
    \end{itemize}
\end{tcolorbox}

\subsection{Synthèse / À retenir}

\begin{itemize}
    \item \textbf{Réutilisabilité} : Les génériques permettent d'écrire un algorithme ou une structure de données \textbf{une seule fois} pour un nombre illimité de types.
    \item \textbf{Sécurité à la compilation} : Contrairement à l'ancienne approche \texttt{object} + cast, les erreurs de type sont détectées par le compilateur et non à l'exécution.
    \item \textbf{Performance} : Les génériques éliminent le boxing/unboxing pour les types valeur, évitant des allocations mémoire superflues sur le tas.
    \item \textbf{Contraintes} : Utilisez \texttt{where} pour restreindre les types autorisés et débloquer l'accès aux méthodes d'interfaces ou de classes de base.
    \item \textbf{Inférence} : Le compilateur déduit automatiquement le type \texttt{T} à partir des arguments. L'écriture explicite n'est nécessaire que lorsqu'il y a ambiguïté.
\end{itemize}
